\section{Hardware results}\label{sec:hardware-results}

\hw{this section is a skeleton: each slot names the number or figure it
waits for; phase A fills \S\ref{sec:hw-map}, phase B
\S\ref{sec:hw-variability}, phase C \S\ref{sec:hw-fewshot} and
\S\ref{sec:hw-same}}

\subsection{The map of robot \#1}\label{sec:hw-map}

Fig.~\ref{fig:map} shows \(u_1\) and \(\omega\) over the \((L,R)\) grid
with \hw{n} replicates per cell, measured over \hw{D} days. Grouped
cross-validation (Table~\ref{tab:hw-cv}) selects the \hw{lin / int /
cor / bat} variant with an error of \hw{x} cm/s in \(u_1\) and \hw{y}
rad/s in \(\omega\), against \hw{x, y} for the best polynomial, \hw{x,
y} for the Gaussian process and \hw{x, y} for the neural network. The
fitted dead zones are \hw{p0L, p0R} and the exponents \hw{gammaL,
gammaR}; the interaction term \(f_L f_R\) \hw{is / is not} supported by
the data. At the top of the grid the speed \hw{stays monotone / reverses
at power P}. The 40 kHz PWM mode \hw{changes only the gain by G \% /
changes the dead zone and exponent}. The mean centre-of-rotation offset
is \(\bar d = \) \hw{d} mm (bootstrap 95 \% interval \hw{lo--hi}); the
step responses \hw{do / do not} separate the clockwise and
counter-clockwise offsets. The lag is \(\tau = \) \hw{tau} s with a dead
time of \hw{t0} s, the heading diffusion \(D_r = \) \hw{Dr} rad\(^2\)/s
\hw{Brownian / OU}. Parameter standard deviations are within \hw{factor}
of the Cramér--Rao bound. Over three days the map drifted by \hw{drift}
in gain and \hw{drift} in dead zone, and the \texttt{bat} term accounts
for \hw{share} of it.

\begin{figure}[t]
\centering
\fbox{\parbox[c][5cm][c]{0.95\columnwidth}{\centering \hw{Fig.~1: \(u_1\) and \(\omega\) over the grid, replicates, fitted map, top-of-grid check}}}
\caption{Robot \#1's forward map. \hw{caption}}
\label{fig:map}
\end{figure}

\begin{table}[t]
\caption{Grouped cross-validation on robot \#1 (RMSE on held-out cells). \hw{fill}}
\label{tab:hw-cv}
\centering\footnotesize
\begin{tabular}{lcc}
\toprule
model & \(u_1\) [cm/s] & \(\omega\) [rad/s] \\
\midrule
\texttt{lin} & -- & -- \\
\texttt{int} & -- & -- \\
\texttt{cor} & -- & -- \\
\texttt{bat} & -- & -- \\
polynomial (plane / quadratic / cubic) & -- & -- \\
Gaussian process & -- & -- \\
neural network & -- & -- \\
\bottomrule
\end{tabular}
\end{table}

\subsection{How different are \hw{N} robots}\label{sec:hw-variability}

Fig.~\ref{fig:variance} and Table~\ref{tab:hw-variance} give the robot
/ day / segment standard deviations per output over \hw{N} robots and
\hw{D} days. The robot component is \hw{x} cm/s in \(u_1\) and \hw{y}
rad/s in \(\omega\), \hw{r} times the segment noise and \hw{r} times
the day component; the fractions are \hw{robot / day / segment}. The
between-robot differences are \hw{dominated by gain / by dead zone /
spread across parameters} (Fig.~\ref{fig:variance}, right), and the
fitted population prior has \hw{SDs per parameter}.

\begin{figure}[t]
\centering
\fbox{\parbox[c][5cm][c]{0.95\columnwidth}{\centering \hw{Fig.~2: variance components per output; per-parameter spread of the fitted maps; the population prior}}}
\caption{Variability across \hw{N} robots and \hw{D} days. \hw{caption}}
\label{fig:variance}
\end{figure}

\begin{table}[t]
\caption{Measured robot, day and segment standard deviations (compare Table~\ref{tab:1}). \hw{fill}}
\label{tab:hw-variance}
\centering\footnotesize
\begin{tabular}{lcccc}
\toprule
output & robot SD & day SD & segment SD & fractions \\
\midrule
\(u_1\) [m/s] & -- & -- & -- & -- \\
\(u_2\) [m/s] & -- & -- & -- & -- \\
\(\omega\) [rad/s] & -- & -- & -- & -- \\
\bottomrule
\end{tabular}
\end{table}

\subsection{Calibrating a new robot}\label{sec:hw-fewshot}

Fig.~\ref{fig:fewshot} and Table~\ref{tab:hw-fewshot} give the
leave-one-robot-out error against the number of segments \(k\), scored
on each robot's measured segments outside the chosen subset. With a
prior from \hw{S = N-1} robots and \hw{k} designed segments (about \hw{t}
minutes of driving) the error reaches \hw{x} noise units, \hw{r} times
better than fitting from scratch at the same budget and \hw{r} times
better than the master/slave standardisation, which plateaus at \hw{x}.
The Klingner-style black-box search reaches \hw{x} after \hw{k}
segments. A prior from one robot \hw{is / is not} enough beyond \(k =
\) \hw{k}. The 95 \% predictive intervals cover the held-out segments
\hw{c} of the time. \hw{if any robot falls outside the model family,
say which and how it showed}

\begin{figure}[t]
\centering
\fbox{\parbox[c][5cm][c]{0.95\columnwidth}{\centering \hw{Fig.~3: error vs \(k\), random vs designed, priors from 1 / 2 / 4 robots, baselines, coverage}}}
\caption{Few-shot calibration of held-out robots. \hw{caption}}
\label{fig:fewshot}
\end{figure}

\begin{table*}[t]
\caption{Measured few-shot error against \(k\), median over held-out robots (compare Table~\ref{tab:2}). \hw{fill}}
\label{tab:hw-fewshot}
\centering\footnotesize
\begin{tabular}{lcccccccc}
\toprule
\(k\) & scratch & black-box & ref\_gain & hier S=1 & hier S=2 & hier S=4 & hier S=1, dopt & hier S=4, dopt \\
\midrule
2 & -- & -- & -- & -- & -- & -- & -- & -- \\
4 & -- & -- & -- & -- & -- & -- & -- & -- \\
8 & -- & -- & -- & -- & -- & -- & -- & -- \\
16 & -- & -- & -- & -- & -- & -- & -- & -- \\
32 & -- & -- & -- & -- & -- & -- & -- & -- \\
64 & -- & -- & -- & -- & -- & -- & -- & -- \\
\bottomrule
\end{tabular}
\end{table*}

\subsection{Same command, same motion}\label{sec:hw-same}

Fig.~\ref{fig:same} shows the achieved \((v,\omega)\) when every robot
is commanded the same targets through its own inverse map, on the
intersection of the feasible sets. Before calibration (one shared map)
the across-robot spread is \hw{x} cm/s and \hw{y} rad/s; after
per-robot calibration it is \hw{x} cm/s and \hw{y} rad/s, \hw{r} times
the segment noise. \hw{whether the residual spread is drift, model
error or noise}

\begin{figure}[t]
\centering
\fbox{\parbox[c][5cm][c]{0.95\columnwidth}{\centering \hw{Fig.~4: achieved vs commanded \((v,\omega)\) per robot, before and after calibration}}}
\caption{Same commanded velocity on all robots. \hw{caption}}
\label{fig:same}
\end{figure}
